\documentclass{article}
\usepackage{graphicx} % Required for inserting images

\title{Un approccio semplificato alla rappresentanza    
(impiegando le Community di Whatsapp)}


\author{Lorenzo Temussi}
\date{November 2026}

\begin{document}


\maketitle


\section{Premessa}
Nella mia breve e faticata carriera di Rappresentante della Paritetica (un organo composto di studenti e docenti che si occupa di fornire un ponte agevolato per la comunicazione fra i due), ho sviluppato un dispositivo semplificato per svolgere i miei oneri istituzionali, e incrementare la qol degli studenti che rappresento con l'impiego del semplicissimo tool della Community di Whatsapp.

Visto poi che per me è un piacere ed un onore condividere i frutti della mia esperienza con i miei colleghi degli altri dipartimenti, ho vergato il documento qui presente per aprire i vantaggi di questo approccio a chiunque sia interessato.

\section{Processo}
Non lasciatevi intimidire dalla lunghezza del processo, richiede poco tempo reale e nessuno strumento che non sia già a vostra disposizione.

\subsection{Designazione dell'Admin}

Il primo passo indispensabile è designare una persona come responsabile o admin della community. Io raccomando di eleggere a tale ruolo una figura di rappresentante già riconosciuta da tutti, io ero già Rappresentante della Paritetica quando ho assunto questo ruolo, la ragione vi verrà spiegata presto.

\subsection{Nomina dell'Admin}
Il secondo passo è conferire a questa persona i diritti di admin di ogni gruppo del dipartimento. Questo è il passaggio più noioso perché richiede di contattare l'admin precedente e chiedere il passaggio dei diritti di admin (funzione "rendi admin", per quanto è un processo una tantum, che non dovrà essere ripetuto. Inoltre molti gruppi sono amministrati dalla stessa persona, per cui basterà fargli una lista dei gruppi in questione e quando avrà 10 minuti morti potrà svolgere il procedimento senza inghippi. Ovviamente l'admin dovrà essere un partecipante in tutti i gruppi per essere adminato.

In casi estremi in cui l'admin di un gruppo non sia disponibile, si raccomanda di creare un nuovo gruppo omonimo, invitare i colleghi nel gruppo non responsivo a trasferirsi su quello appena creato tramite link d'invito (opzione presente dal menu "aggiungi membro" del nuovo gruppo), e procedere come se questo fosse il gruppo originale.

\subsection{Creazione della Community}
Il terzo passo, affidato al nuovo admin universale, è recarsi sul menu principale di Whatsapp e utilizzare il pulsante "crea nuovo", ovvero "+", avendo cura di specificare "Nuova Community". In fase di creazione sarà richiesto di aggiungere gruppi amministrati dall'utente, qui è possibile aggiungerli subito tutti o solo una parte, e aggiungere i rimanenti a miglior convenienza in seguito.

Questo processo richiede circa 5 secondi.

\subsection{Suddivisione in canali}
Per il quarto passo, i gruppi nella community, che chiameremo canali, da adesso in avanti, potrebbero essere in difetto o in eccesso rispetto al dovuto, fate in modo che rispondano meglio alle reali necessità del corso, condivido il mio processo, ma ognuno dovrà usare il proprio buonsenso:

- Un gruppo per ogni anno del corso di studi (1° anno, 2°anno, etc...)
- Un gruppo Fuori Corso e un gruppo Magistrali (almeno)

Alternativamente

- Un gruppo per ogni anno di immatricolazione (2025, 2026, 2027, etc...)
- Un gruppo solo per i gruppi "stagionati", per evitare di avere 30 gruppi sulla community e tornare punto e a capo.

In aggiunta:

- Un gruppo Compravendita Libri/Appunti
- Un gruppo Chiacchiere, Meme o Discorsi Leggeri per la socializzazione anche prolissa, in modo da evitare di offuscare comunicazioni di interesse pubblico sui gruppi principali.
- Un gruppo per organizzare attività extracurricolari, che siano uscite al centro, partecipazione a eventi pubblici o simili. Qui dentro verranno anche condivisi link a gruppi esterni nel caso sia necessaria una organizzazione più dettagliata.

\section{Mantenimento}

Ora che avete la vostra community proprio come quella di informatica, potete fare un passo in più e applicare ulteriori buone pratiche per mantenerla nello stato migliore possibile. Questi consigli sono rivolti all'admin, ma non ristretti all'admin.

\subsection{Contratto d'uso}

Designate con chiarezza e stringatezza il caso d'uso di ogni canale nella descrizione. Un buon inizio potrebbe essere: "Questo canale è ad uso degli studenti UNIPG per organizzarsi e socializzare".

Potete anche aggiungere un link utile in descrizione, ad esempio quello del gruppo dell'UDU, oppure ad un gruppo telegram esterno con delle risorse d'interesse pubblico, o ancora a delle cartelle di appunti per il vostro corso di studi (a informatica gli appunti di tutte la materie sono stati messi a disposizione degli altri studenti gratuitamente, ma è un po un caso a parte).

\subsection {Pulizia}

Quando un utente pubblica contenuti inadatti a canale, abbiate la bontà di riprenderlo (anche privatamente se preferite) facendo riferimento al caso d'uso in descrizione e al feedback ricevuto dagli altri utenti.

Se un utente persiste nella sua follia, potete cacciarlo da tutti i gruppi con un semplice clic. Abbiate anche la pazienza di eleminare eventuali contenuti inappropriati, se non siete sicuri che siano inappropriati indite una poll per capire l'opinione pubblica, ma ricordate che l'autorità è, in ultimo, in mano a voi.

\subsection{Convivialità}

Avrete forse notato che in taluni gruppi il dialogo fra i membri non nasce spontaneo. Può esserci un numero di ragioni, che non ci interessano.

Una volta eliminati eventuali elementi ostili dal gruppo (bot che vende i biglietti per il concerto di Bugs Bunny a Milano, mentecatto che assilla le studentesse in privato et simili), cercate di creare un po' di investimento da parte dei membri. Quando qualcuno fa una domanda cercate effettivamente di rispondere correttamente, fate osservazioni sulla lezione del giorno o anche sulla partita della sera prima. Non parlate di religione o di politica.


Vedrete che dopo un po' si creerà un piccolo gruppo di persone che tratta la chat come attiva e utile a qualcosa, a quel punto sarà a tutti gli effetti sostenibile in autonomia. Pensate all'attività iniziale come quel giochino che si fa ai gatti per fargli bere l'acqua che metti un dito nella ciotola e fai biribiribiri finché non non ha ricevuto il messaggio che è roba buona.

\section{Vantaggi}

I vantaggi non si contano, ma proviamo a contarli comunque

\subsection{Pure Metriche}

Gli studenti di Informatica, presentati con un questionario di gradimento hanno votato quasi all'unanimità una marcata preferenza per il nuovo approccio (90\% a favore circa).

\subsection{Risparmio di Tempo}

Gli studenti di tutti gli anni possono contattare il rappresentante di turno in pochi secondi, che a sua volta può velocemente ridirigerli alle risorse che richiedono ("boh nlo so chiedi sul gruppo").

Il Rappresentante ha sotto mano una serie di validissimi strumenti per comunicazioni one-to-many e viceversa: può informare centinaia di studenti con un solo post in bacheca e raccoglierne le opinioni con una singola poll ben strutturata.

Dalla mia esperienza, fintanto che mantengo i gruppi privi di spam e appena moderati, i colleghi fanno praticamente il mio lavoro al mio posto. Anche quando lancio una qualche investigazione più approfondita (ad esempio su una sessione d'esame andata a rotoli), la community mi permette di connettermi più facilmente alle persone d'interesse.

In aggiunta un numero sufficientemente alto di studenti può concordare in autonomia quale sia una problematica reale e notificarle al Rappresentante, che non deve fare altro che annotarla.

\subsection{Produzione di metriche accurate}
Come sopra, i poll "globali" offrono dati quantificabili altrimenti pressoché impossibili da rilevare.


\subsection{Streamlining degli accessi}
Uno studente del primo anno che entra nella community ha subito acesso a tutti i gruppi, senza dover chiedere link di sorta a chicchessia.

\section{Ringraziamenti e saluti}
Ringrazio i colleghi per aver letto il presente documento e li invito a metterlo in pratica per diminuire il volume di lavoro (particolarmente il lavoro inutile) legato all'attività rappresentativa e studentesca.

Mi lascio a disposizione per ogni dubbio o curiosità, mi trovate facilmente nella community di Informatica.

Ringrazio inoltre Alessandro, Raul, e Francesco che hanno dimostrato grande spirito prosociale e mi hanno assistito ed ispirato a sviluppare strumenti per agevolare le parti più scomode della vita universitaria.
\end{document}
